\begin{pastexam}[label={px:ntu-109-comm-dsa-34}]{臺大電信所 109 學年度（B）第 34 題}
（單選；每題 2 分，答錯不倒扣）

Consider the following list of 9 activities with their start time $s$ and end time $t$, what's the
maximum number of activities that can be performed by a single person? Assuming that a person can
only work on a single activity at a time.

\begin{center}
\begin{tabular}{|c|c|c|c|c|c|c|c|c|c|}\hline
\textbf{index} & 1 & 2 & 3 & 4 & 5 & 6 & 7 & 8 & 9\\\hline
$s$ & 2 & 3 & 1 & 10 & 14 & 6 & 13 & 7 & 19\\\hline
$t$ & 5 & 9 & 12 & 16 & 17 & 15 & 18 & 11 & 20\\\hline
\end{tabular}
\end{center}
(A)~2 \qquad (B)~3 \qquad (C)~4 \qquad (D)~5
\end{pastexam}

\begin{pastsoln}
答案是 (C)。

這是活動選擇問題。 貪心法：依結束時間由早到晚排，每次選「開始時間不早於上一個選中活動的
結束時間」的第一個活動。 選最早結束的活動一定有某個最佳解包含它，所以貪心的結果是最佳解
（資結書定理~19.1）。

依結束時間排序並逐一檢查（「上一個結束」是目前已選活動中最後結束的時間）：
\begin{center}
\begin{tabular}{cccccl}
\toprule
活動 & $s$ & $t$ & 上一個結束 & 選？ & 理由 \\
\midrule
1 & 2 & 5 & -- & 選 & 第一個 \\
2 & 3 & 9 & 5 & 不選 & $3 < 5$ \\
8 & 7 & 11 & 5 & 選 & $7 \ge 5$ \\
3 & 1 & 12 & 11 & 不選 & $1 < 11$ \\
6 & 6 & 15 & 11 & 不選 & $6 < 11$ \\
4 & 10 & 16 & 11 & 不選 & $10 < 11$ \\
5 & 14 & 17 & 11 & 選 & $14 \ge 11$ \\
7 & 13 & 18 & 17 & 不選 & $13 < 17$ \\
9 & 19 & 20 & 17 & 選 & $19 \ge 17$ \\
\bottomrule
\end{tabular}
\end{center}
選出活動 $1, 8, 5, 9$，時間是 $[2, 5]$、$[7, 11]$、$[14, 17]$、$[19, 20]$，兩兩不重疊，共
$4$ 個。 貪心法是最佳的，所以最多 $4$ 個，選 (C)。 這裡沒有哪個活動的開始時間剛好等於另一個的
結束時間，所以端點算不算重疊不影響答案。
\end{pastsoln}
